\subsection{Non-injective projection}
  \label{subsec:non-injective-projection}

  Observable outcomes are defined through a projection
  \[
    \Pi : \Omega \rightarrow O_A \times O_B ,
  \]
  where $O_A$ and $O_B$ denote the sets of possible outcomes for measurements performed
  in regions $A$ and $B$, respectively.

  The projection $\Pi$ is assumed to be non-injective.
  Distinct underlying configurations may therefore correspond to the same pair of
  observable outcomes $(a,b)$.
  Such non-injectivity arises naturally when the effective description $O_A \times O_B$
  has a finite capacity to encode distinctions present in $\Omega$.

  Observable states are thus identified with equivalence classes
  \[
    [\omega] = \{ \omega' \in \Omega \mid \Pi(\omega') = \Pi(\omega) \}.
  \]
